\documentclass[8pt,a4paper,twocolumn]{ltjsarticle}

\usepackage[
a4paper,
margin=16.5mm
]{geometry}

\usepackage{newtxtext,newtxmath}
\usepackage[hidelinks]{hyperref}
\usepackage{url}
\usepackage{titlesec}
\usepackage{microtype}
\usepackage{cite}
\usepackage{booktabs}
\usepackage{array}
\usepackage{enumitem}

% 2段組の段間
\setlength{\columnsep}{6mm}

% 標準に近い段落設定
\setlength{\parindent}{1em}
\setlength{\parskip}{0pt}
\setlength{\emergencystretch}{1em}

\setlist[itemize]{
leftmargin=1.5em,
itemsep=1pt,
topsep=2pt,
parsep=0pt
}

% 見出しの前後を自然な間隔にする
\titlespacing*{\section}{0pt}{5pt}{2pt}
\titleformat{\section}
{\bfseries\normalsize}
{\thesection.}
{0.5em}
{}

\renewcommand{\refname}{参考文献}

\begin{document}

\twocolumn[
\begin{center}
{\large\bfseries
大規模道路データ応用としての舗装点検自動化システム「ロードタイガー」
\par}

\vspace{1pt}

\textbf{氏名：}杉山月渚
\qquad
\textbf{学籍番号：}48-256422
\end{center}

\vspace{0.5mm}
]

\section{対象サービスと背景}

ロードタイガー（ROAD TIGER）は、NEXCO中日本と中日本ハイウェイ・エンジニアリング東京が開発・運用する走行型の路面性状測定システムである。カメラ、レーザー、加速度計、測位装置などを車両に搭載し、交通流の中を走行しながら、舗装の代表的な劣化指標であるひび割れ、わだち掘れ、平坦性を連続的に測定する。1982年に1号車が開発され、2023年に公表された新型は7号車であり、同年11月から東名高速道路などで実測を開始した実用例である\cite{nexco2023}。

\section{データ取得・解析の仕組み}

新型車の中核は、左右2台のラインスキャンカメラで構成される3Dステレオカメラである。青色LEDを路面に照射して昼夜で明るさをそろえ、左右画像の視差から路面の高さ画像と横断形状を算出する。これにより、1台のカメラ系でひび割れ画像とわだち掘れの横断プロファイルを取得できる。平坦性は、車両左右の非接触レーザー変位計と加速度計を用い、車体振動を補正しながら3mシグマおよびIRI（国際ラフネス指数）として評価する。GNSSとIMUは位置、姿勢、縦横断勾配を補完する\cite{catalog2025,roadtigerpdf}。

収集後は、撮影画像からひび割れを抽出してひび割れ率を計算し、高さ情報からわだち深さ、レーザー計測から平坦性を算出する。結果をキロポストや位置情報と結び付ければ、区間別の状態評価、過年度データとの比較、補修候補の抽出、工事後の効果検証まで同じデータ系列で行える。新型では解析処理の大部分が自動化され、今後はレーンマークの剥がれ、骨材飛散、局部沈下、段差などへの機能拡張も想定されている\cite{nexco2023,catalog2025}。

\section{規模を示す定量的エビデンス}

測定速度は従来の60--100 km/hから30--120 km/hへ拡大し、昼夜を問わず測定できる。道路表面は横断方向4.4 m、進行方向1.25 mm間隔で撮像し、幅1 mmのひび割れを確認可能である\cite{nexco2023,catalog2025}。1.25 mm間隔は1 m当たり800本、1 km当たり80万本の縦断撮像ラインに相当する。120 km/h（毎秒約33.3 m）では、仕様値から単純計算して\textbf{毎秒約2.67万ライン}を取得する密度となる。

NEXCO中日本の営業延長は2026年4月時点で2,208 km、利用台数は2024年度実績で約205万台/日である\cite{nexcoarea}。仮に営業延長2,208 kmを1車線分だけ1回、1.25 mm間隔で取得しても、進行方向の撮像ラインは約17.7億本となる。実際の高速道路は上下線と複数車線を持ち、画像、高さ、加速度、位置情報を同時に記録するため、ロードタイガーは典型的な大規模時空間データサービスといえる。

\section{業務上の効果}

第一に、交通規制を極力増やさず高速で測定できる。通常の交通速度に近い範囲で走行できるため、低速点検車と一般車両の速度差を小さくし、後続車の急減速や点検員の路上作業リスクを抑えられる。速度範囲が30 km/hまで広がったことで、渋滞時や道路条件に応じた柔軟な測定も可能になった\cite{nexco2023}。

第二に、昼間測定と自動解析によって生産性が向上した。旧型はひび割れの陰影を得るため夜間中心であったが、ステレオカメラと青色LEDにより昼間も同品質の画像を取得できる。機器の集約と解析自動化により、測定体制は3名から2名へ33\%削減され、車体の小型化で普通免許による運転も可能となった。解析時間の短縮、夜間勤務の削減、車両・人員コストの低減を同時に期待できる\cite{nexco2023}。

第三に、データ品質を継続的に管理している。ロードタイガーの測定性能は土木研究センターの年1回の性能確認試験で確認され、測定機器メーカーによる年1回の点検、自社による1--2カ月ごとの性能確認、乗車前の動作確認が行われる\cite{nexcosafety}。大規模データは量だけでなく、年度間で比較できる再現性と校正が重要であり、この運用は意思決定に使えるデータの信頼性を支えている。

\section{社会的意義}

最大の意義は、事故や走行性悪化の前に補修する予防保全を支えることである。ひび割れやわだちは排水性、操縦安定性、乗り心地を低下させ、放置すれば損傷が表層から路盤へ進行する。国土交通省は、ひび割れ率が40\%を超えた状態を放置すると、雨水の影響で路盤から細粒分が流出し、表層だけを繰り返し補修しても早期に再劣化する可能性を指摘している\cite{xroad2022}。高密度な定期データで小さな変化を捉えれば、損傷が深部に達する前に対処し、ライフサイクルコストを抑えられる。

社会全体でも舗装修繕の需要は大きい。国土交通省の2020年調査では、地方公共団体管理道路のうち修繕段階にある舗装は約5.5万km、そのうち路盤以下の損傷が想定される区間は約2万kmであった\cite{mlit2020}。限られた予算と技術者で全区間を同時に直すことはできないため、客観データに基づく優先順位付けが不可欠である。ロードタイガーのようなシステムは、担当者ごとの目視差を減らし、経年変化、交通量、補修履歴を組み合わせて、危険度と費用対効果の高い区間から投資する根拠を与える。

さらに、測定結果を位置付きで蓄積することは、補修前後の性能比較、材料・工法の耐久性評価、予算執行の説明責任にも役立つ。規制時間の短縮は道路利用者の渋滞損失を抑え、高速走行測定は作業員が車外に出る機会を減らす。したがって本システムは、交通安全、労働安全、維持管理の省人化、公共投資の合理化を同時に支える社会基盤サービスである。

\section{課題と今後の展開}

一方で、雨天時や路面が乾燥していない場合は計測できず、照明条件、汚れ、補修跡、影などが画像解析に影響する可能性がある\cite{catalog2025}。自動判定の誤検出・見逃しを管理するため、学習データの更新、定期校正、専門技術者による確認は残る。また、高密度画像を長期間保存する費用、旧型・他社装置とのデータ互換性、評価指標や帳票形式の標準化も必要である。

専用車両は高性能である反面、導入・維持費が大きく、狭い生活道路や小規模自治体へそのまま展開することは難しい。将来は、専用車による高精度定期測定と、一般車両・スマートフォン・安価なセンサによる高頻度スクリーニングを組み合わせ、共通データ基盤で管理する形が有効である。国土交通省が進めるxROADのような基盤と接続し、点検、診断、設計、施工、品質管理を一貫してデジタル化できれば、ロードタイガーの価値は単独の測定車から全国的な舗装アセットマネジメントへ拡張する\cite{xroad2022}。

\section{結論}

ロードタイガーは、最高120 km/h、4.4 m幅、1.25 mm間隔という高密度測定と、自動解析・位置管理を組み合わせた実在の大規模データ応用である。2,208 kmの高速道路網と1日約205万台の交通を支える環境で、損傷の早期発見、補修優先順位の客観化、測定の省人化、安全な作業、長期的な維持費削減を実現する点に大きな社会的意義がある。

{\fontsize{6}{6.5}\selectfont

\setlength{\itemsep}{0pt}
\setlength{\parskip}{0pt}

\begin{thebibliography}{9}

\bibitem{nexco2023}
NEXCO中日本・中日本ハイウェイ・エンジニアリング東京，
``路面性状測定車「ロードタイガー」をフルモデルチェンジ，''
2023年10月6日．
\url{https://www.c-nexco.co.jp/corporate/pressroom/news_release/5808.html}

\bibitem{catalog2025}
高速道路調査会，
``道路インフラ測定システム ユニオンツール ロードタイガー，''
新技術情報，2025年10月15日更新．
\url{https://catalog.express-highway.or.jp/products/p10015017.html}

\bibitem{roadtigerpdf}
中日本ハイウェイ・エンジニアリング東京，
``路面性状測定車 ROAD TIGER，''
2023年10月．
\url{https://www.c-nexco-het.jp/wp-content/uploads/2023/10/58b50ec91452f1e46a9eeb5703c1f6b1.pdf}

\bibitem{nexcosafety}
NEXCO中日本，
``高速で測定できる、路面性状測定車「ロードタイガー」，''
安全性向上の取組み．
\url{https://www.c-nexco.co.jp/corporate/safety/torikumi/torikumi/vol04/}

\bibitem{nexcoarea}
NEXCO中日本，
``会社概要：事業データ，''
営業延長は2026年4月1日現在、利用台数は2024年度実績．
\url{https://www.c-nexco.co.jp/corporate/employment/new-graduates/company/}

\bibitem{xroad2022}
国土交通省道路局，
``xROADを活用した次世代の舗装マネジメント，''
2022年．
\url{https://www.mlit.go.jp/road/ir/ir-council/dourogijutsu/pdf09/04.pdf}

\bibitem{mlit2020}
国土交通省，
``予防保全によるメンテナンスへの転換について，''
2020年．
\url{https://www.mlit.go.jp/policy/shingikai/content/001375673.pdf}

\end{thebibliography}
}

\end{document}
